% ============================================================
% Chapter: Local LLM Infrastructure with Ollama
% ============================================================
% Preamble requirements (add to main document _preamble.tex):
%   \usepackage{listings}
%   \usepackage{xcolor}
%   \usepackage{booktabs}
%   \usepackage{amsmath}
%   \usepackage{hyperref}
%   \usepackage{tcolorbox}
%   \tcbuselibrary{skins, breakable}
%
% \lstdefinestyle{bashstyle}{...}   % see preamble template
% \lstdefinestyle{psstyle}{...}     % PowerShell variant
% \newtcolorbox{warningbox}{...}
% \newtcolorbox{infobox}{...}
% ============================================================

\chapter{Local LLM Infrastructure with Ollama}
\label{ch:ollama-infrastructure}

\section{Introduction}

\vs
\noindent
\textbf{For students:} this chapter has two parts relevant to you.
Section~\ref{sec:apikeys} explains how to configure the API keys
required to run the agent programs in
\texttt{resources/unit\_7/FOO/}. The Ollama CLI reference in
Section~\ref{sec:install} covers the commands you will use during
the LLM unit exercises. You do not need to read the remainder of
this chapter to complete the data challenge.
\vs

\vs
\noindent
\textbf{For instructors and program staff:} this chapter documents
the full configuration of the two-machine LLM serving infrastructure
used by the \ProjectName{} program, covering the Dell
Precision 7770 application host and the NVIDIA Spark inference
server. It includes Ollama installation, \texttt{systemd} service
configuration, network access control, memory budget management,
single-machine and two-machine deployment modes, connectivity
testing, and capacity planning.
\vs

\section{Hardware Inventory}
\label{sec:hardware}

The two machines participating in this setup have fundamentally different
memory architectures, which governs model selection and agent capacity.

\subsection{Dell Precision 7770 (Application Host)}

\begin{itemize}
  \item \textbf{CPU RAM:} 128\,GB (112\,GB available after OS overhead)
  \item \textbf{GPU:} NVIDIA RTX A5500 Laptop GPU, 16\,GB VRAM (discrete)
  \item \textbf{Architecture:} x86\_64
  \item \textbf{OS:} Windows~11 with WSL2 Ubuntu~24.04
  \item \textbf{Role:} Runs a FastAPI application,
        PostgreSQL~16, and the  agent
        orchestration layer. Hosts a local Ollama instance for
        single-machine development mode.
\end{itemize}

On the Dell, the RTX~A5500 VRAM (16\,GB) is \emph{separate} from system
RAM. Ollama loads model weights into VRAM first and spills to system RAM
only when the model exceeds VRAM capacity. Inference degrades
significantly once spilling begins: tokens per second can drop by an
order of magnitude for spilled portions of the model.

\subsection{NVIDIA Spark (LLM Server)}

\begin{itemize}
  \item \textbf{Unified memory:} 128\,GB physical; approximately 120\,GB
        visible to Ollama (110\,GB conservatively available)
  \item \textbf{GPU:} NVIDIA GB10 SoC
  \item \textbf{Architecture:} ARM64 (AArch64)
  \item \textbf{OS:} Ubuntu~24.04 (native, no virtualisation)
  \item \textbf{Role:} Dedicated LLM inference server. CPU and GPU share
        the same memory pool, so the entire 110\,GB is available for
        model weights and KV cache simultaneously.
\end{itemize}

\begin{warningbox}
The GB10 uses a \emph{unified memory} architecture. Unlike a discrete
GPU, there is no separate VRAM budget to track. The 110\,GB pool covers
weights, KV cache, and OS overhead together. This makes the Spark well
suited to serving multiple large models simultaneously (provided context
windows are bounded), but it also means that vLLM's PagedAttention,
which is optimised for discrete high-VRAM GPUs, provides fewer benefits
here than Ollama's native unified-memory handling.
\end{warningbox}

\section{Network Topology}
\label{sec:network}

Both machines reside on the UTSA \texttt{utsarr.net} subnet. At the time
of configuration they held the following addresses:

\begin{center}
\begin{tabular}{llll}
\toprule
Machine & Interface & IP Address & Subnet \\
\midrule
Spark   & \texttt{enP7s7} & \texttt{10.2.14.24}  & \texttt{/23} \\
Dell    & Ethernet~8      & \texttt{10.2.14.143} & \texttt{/23} \\
\midrule
Spark   & \texttt{enP7s7} MAC & \multicolumn{2}{l}{\texttt{4c:bb:47:2c:74:eb}} \\
Dell    & Ethernet~8 MAC      & \multicolumn{2}{l}{\texttt{ac:91:a1:7d:06:ff}} \\
\bottomrule
\end{tabular}
\end{center}

\subsection{Switch vs.\ Router}

A Netgear ProSAFE GS105 unmanaged 5-port switch connects the two
machines. An \emph{unmanaged switch} operates at OSI Layer~2: it
forwards Ethernet frames between ports based on MAC addresses but has no
concept of IP subnets, routing, or traffic restriction. A
\emph{router} operates at Layer~3 and is required to separate subnets or
enforce IP-level traffic policies.

Because both machines are on the same \texttt{/23} subnet managed by the
university network, ARP entries are not visible between hosts; traffic
crosses at least one router hop, and ARP does not traverse routers. This
is expected behaviour on a routed network and should not be interpreted
as a connectivity error.

\subsection{MAC Address Filtering}

MAC-based firewall rules (via \texttt{iptables}) are only meaningful on
the same Layer~2 segment. They can be spoofed trivially and should be
treated as an informational layer rather than a security boundary. IP
firewall rules (\texttt{ufw}) are the appropriate primary control for
restricting Ollama access to known client addresses.

\begin{lstlisting}[style=bashstyle, caption={Restrict Ollama to the Dell's IP only}]
# On the Spark: allow only the Dell, deny all other sources
sudo ufw allow from 10.2.14.143 to any port 11434
sudo ufw deny 11434
\end{lstlisting}

\subsection{Identifying Network Interfaces}

Use the following commands to enumerate interfaces and confirm addresses
before configuring firewall rules.

\begin{lstlisting}[style=bashstyle, caption={Linux: list interfaces and addresses}]
ip a
# Look for interfaces with state UP and a valid inet address.
# The link/ether line on each interface gives the MAC address.
\end{lstlisting}

\begin{lstlisting}[style=psstyle, caption={Windows: list interfaces and addresses}]
ipconfig /all
# Look for active Ethernet adapters with an IPv4 Address.
# The Physical Address field gives the MAC (dash-separated format).
\end{lstlisting}

\begin{infobox}
\textbf{MAC address formats.} Linux tools (\texttt{ip~a},
\texttt{arp}) report MACs in lowercase colon format:
\texttt{ac:91:a1:7d:06:ff}. Windows \texttt{ipconfig~{/}all} reports
them in uppercase dash format: \texttt{AC-91-A1-7D-06-FF}. Both
represent the same address. Configuration files and scripts should
normalise to a single format before comparison.
\end{infobox}

\section{Installing Ollama}
\label{sec:install}

\subsection{Linux (Spark)}

The official installer script detects the architecture automatically and
selects the appropriate binary (ARM64 in this case).

\begin{lstlisting}[style=bashstyle, caption={Install Ollama on Ubuntu}]
curl -fsSL https://ollama.com/install.sh | sh
\end{lstlisting}

The installer creates the \texttt{ollama} user and group, adds the
current user to the \texttt{ollama} group, installs the binary to
\texttt{/usr/local/bin/ollama}, and registers a \texttt{systemd} service.
Confirm the installation:

\begin{lstlisting}[style=bashstyle]
ollama --version
systemctl status ollama
\end{lstlisting}

\subsection{Windows (Dell)}

Download the Windows installer from \url{https://ollama.com} and run it.
Ollama registers itself in the system tray and starts automatically. The
API listens on \texttt{127.0.0.1:11434} by default.

\subsection{Conda Environments}

The Spark ships with Miniconda, which activates a \texttt{(base)}
environment by default. Ollama does not require Python and should be
installed at the system level, not inside a Conda environment. Disable
automatic Conda activation to avoid interference with system tooling:

\begin{lstlisting}[style=bashstyle]
conda config --set auto_activate_base false
\end{lstlisting}

For all Python work in this stack, use plain \texttt{venv} environments,
which translate directly to \texttt{systemd} \texttt{ExecStart} paths
and maintain production parity without Conda's dependency overhead.

\section{Configuring Ollama for Network Access}
\label{sec:network-config}

By default Ollama binds only to \texttt{localhost}. For the two-machine
setup the Spark must expose its API to the network. Configure this
through a \texttt{systemd} override so the setting survives service
restarts and reboots.

\begin{lstlisting}[style=bashstyle, caption={Create a systemd override on the Spark}]
sudo systemctl edit ollama
\end{lstlisting}

This opens \texttt{/etc/systemd/system/ollama.service.d/override.conf}.
Add the following:

\begin{lstlisting}[style=bashstyle, caption={Spark Ollama systemd override}]
[Service]
Environment="OLLAMA_HOST=0.0.0.0:11434"
Environment="OLLAMA_MAX_LOADED_MODELS=4"
Environment="OLLAMA_NUM_PARALLEL=4"
Environment="OLLAMA_KEEP_ALIVE=24h"
Environment="OLLAMA_NUM_CTX=8192"
\end{lstlisting}

Apply the changes:

\begin{lstlisting}[style=bashstyle]
sudo systemctl daemon-reload
sudo systemctl restart ollama
systemctl status ollama
\end{lstlisting}

The environment variables are explained in Table~\ref{tab:ollama-env}.

\begin{table}[h]
\centering
\caption{Key Ollama environment variables}
\label{tab:ollama-env}
\begin{tabular}{lp{9cm}}
\toprule
Variable & Effect \\
\midrule
\texttt{OLLAMA\_HOST} &
  Bind address and port. Use \texttt{0.0.0.0:11434} to accept connections
  from any network interface. \\
\texttt{OLLAMA\_MAX\_LOADED\_MODELS} &
  Maximum number of models resident in memory simultaneously. Prevents
  thrashing when \textsc{Max} dispatches concurrent agent calls. \\
\texttt{OLLAMA\_NUM\_PARALLEL} &
  Maximum concurrent inference requests per model. \\
\texttt{OLLAMA\_KEEP\_ALIVE} &
  Duration a model remains loaded after its last request. The default is
  5 minutes; \texttt{24h} keeps models warm across a workday. \\
\texttt{OLLAMA\_NUM\_CTX} &
  Default context window in tokens. \textbf{Critical on the Spark}: without
  this cap, Ollama defaults to 262,144~tokens, which allocates a KV cache
  large enough to exceed physical memory on most models. Set to
  \texttt{8192} as a safe default. \\
\bottomrule
\end{tabular}
\end{table}

\section{Memory Budget and Context Window Management}
\label{sec:memory}

\subsection{The Governing Formula}

Total memory consumed per agent has two independent components:

\begin{equation}
M_{\text{agent}} = M_{\text{weights}} + M_{\text{KV}}
\label{eq:memory}
\end{equation}

\noindent where the weight footprint under Q4 quantisation is
approximated by:

\begin{equation}
M_{\text{weights}} \approx 0.5\,\text{GB} \times B_{\text{params}}
\label{eq:weights}
\end{equation}

\noindent and the KV cache footprint is:

\begin{equation}
M_{\text{KV}} \approx \frac{n_{\text{layers}} \times n_{\text{kv\_heads}}
  \times d_{\text{head}} \times C \times 2 \times 2}{10^9}\;\text{GB}
\label{eq:kvcache}
\end{equation}

\noindent where $C$ is the context length in tokens and the factor of 2
appears twice for the key/value pair and the 2-byte (BF16) dtype.

\begin{warningbox}
\textbf{The context window dominates memory for large models.} A
\texttt{llama3.3:70b} model requires approximately 40\,GB for weights
at Q4. With Ollama's default context of 262,144~tokens, the KV cache
adds another 134\,GB, pushing total allocation to 174\,GB, exceeding
the Spark's 120\,GB and causing the load to fail. With a context of
8,192~tokens, the KV cache drops to 2.5\,GB, giving a total of
42.5\,GB which fits comfortably.
\end{warningbox}

\subsection{The Modelfile Pattern for Large Models}

Ollama does not support runtime context overrides via the
\texttt{--num-ctx} flag. The correct mechanism is a \emph{Modelfile},
which creates a named derivative of the base model with the parameter
baked in.

\subsubsection{Linux (Spark and WSL)}

\begin{lstlisting}[style=bashstyle, caption={Create a context-bounded model variant on Linux}]
cat > /tmp/Modelfile.llama70b << 'EOF'
FROM llama3.3:70b
PARAMETER num_ctx 8192
EOF
ollama create llama3.3:70b-ctx8k -f /tmp/Modelfile.llama70b
ollama run llama3.3:70b-ctx8k "Reply with one word: ready"
\end{lstlisting}

\subsubsection{Windows (PowerShell)}

PowerShell uses a here-string syntax for multi-line literals. The
closing \texttt{"@} must appear at the \emph{absolute start of a line}
with no leading spaces or tabs; this is a hard requirement of the
PowerShell parser.

\begin{lstlisting}[style=psstyle, caption={Create a context-bounded model variant on Windows}]
@"
FROM llama3.3:70b
PARAMETER num_ctx 8192
"@ | Out-File -FilePath "$env:TEMP\Modelfile.llama70b" -Encoding utf8
ollama create llama3.3:70b-ctx8k -f "$env:TEMP\Modelfile.llama70b"
\end{lstlisting}

Apply the same pattern to every model above 20\,B parameters. The
naming convention \texttt{<model>-ctx8k} makes the operational
constraint explicit at the call site.

\subsection{Hardware Agent Capacity}

Given the memory budget formula, the maximum number of agents that can
run simultaneously on a machine is:

\begin{equation}
N_{\text{hw}} = \left\lfloor \frac{M_{\text{available}}}{M_{\text{agent}}} \right\rfloor
\label{eq:nhw}
\end{equation}

Table~\ref{tab:model-registry} gives the confirmed model inventory with
measured status on each platform and approximate memory footprint at
8,192-token context.
{\scriptsize
\begin{table}[h]
\centering
\caption{Confirmed model inventory and memory footprint at 8K context}
\label{tab:model-registry}
\footnotesize
\begin{tabular}{llrrrll}
\toprule
Model & Ollama name & Params & $M_w$ & $M_{\text{KV}}$ &
  $M_{\text{total}}$ & Status \\
 & & (B) & (GB) & (GB) & (GB) & \\
\midrule
nemotron-3-super:120b & \texttt{...-ctx8k} & 120 (MoE) & 65 & 4.0 & 69.0 &
  \checkmark\ both \\
llama3.3:70b          & \texttt{...-ctx8k} & 70         & 40 & 2.5 & 42.5 &
  \checkmark\ both \\
qwen3:32b             & \texttt{...-ctx8k} & 32         & 18 & 1.3 & 19.3 &
  \checkmark\ both \\
mistral-small:22b     & \texttt{mistral-small:22b} & 22 & 13 & 1.0 & 14.0 &
  \checkmark\ both \\
phi4:14b              & \texttt{phi4:14b}  & 14         &  8 & 0.8 &  8.8 &
  \checkmark\ both \\
qwen3:14b             & \texttt{qwen3:14b} & 14         &  8 & 0.8 &  8.8 &
  \checkmark\ both \\
mistral-nemo:12b      & \texttt{mistral-nemo:12b} & 12 &  7 & 0.7 &  7.7 &
  \checkmark\ both \\
mistral:7b            & \texttt{mistral:7b}& 7          &  4 & 0.5 &  4.5 &
  \checkmark\ both \\
deepseek-r1:8b        & \texttt{deepseek-r1:8b} & 8     &  5 & 0.5 &  5.5 &
  \checkmark\ both \\
gemma3:4b             & \texttt{gemma3:4b} & 4          &  2.5 & 0.3 & 2.8 &
  \checkmark\ both \\
phi4-mini:3.8b        & \texttt{phi4-mini:3.8b} & 3.8   &  2.5 & 0.3 & 2.8 &
  \checkmark\ both \\
qwen3:4b              & \texttt{qwen3:4b}  & 4          &  2.5 & 0.3 & 2.8 &
  \checkmark\ Dell; \dag\ Spark \\
gemma3:1b             & \texttt{gemma3:1b} & 1          &  0.8 & 0.1 & 0.9 &
  \checkmark\ both \\
deepseek-v3.1:671b    & ---                & 671 (MoE)  & 404 & 8.0 & 412 &
  $\times$\ too large \\
\bottomrule
\multicolumn{7}{l}{\dag\ ARM64 crash on Spark (Ollama bug); works on x86\_64 Dell.}
\end{tabular}
\end{table}
}
\subsection{Thinking Models and Token Budgets}

Several models in the inventory implement chain-of-thought reasoning
internally, emitting a \emph{thinking block} before their final answer.
These include \texttt{qwen3} (all sizes), \texttt{deepseek-r1}, and
\texttt{nemotron-3-super}. When calling these models via the
OpenAI-compatible endpoint with a small \texttt{max\_tokens} budget, the
budget may be exhausted inside the thinking block, producing either an
HTTP~500 error or an empty reply.

The correct mitigation is \emph{not} to suppress thinking via the
\texttt{/no\_think} system prompt, which can leave the model with nothing
to say outside the block. Instead, set \texttt{max\_tokens} generously
(500 or more) and strip thinking tags from the response:

\begin{lstlisting}[style=bashstyle, caption={Thinking tag patterns to strip from model output}]
# DeepSeek-R1 style:
<think> ... </think>

# Qwen3 / Ollama native style:
Thinking...
  ... reasoning text ...
...done thinking.
\end{lstlisting}

\section{Configuring API Keys}
\label{sec:apikeys}

You will receive an API key for OpenAI and Anthropic during our first session. You can get your own API for free at \href{https://groq.com}{Groq}. You need to create three API key variables following the procedure described below: \texttt{OPENAI\_API\_KEY}, \texttt{ANTHROPIC\_API\_KEY}, \texttt{GROQ\_API\_KEY}. Test the programs located at \texttt{resources/unit\_7/FOO/} in the \ProjectName{} repository: \texttt{agentGroq.py}, \texttt{agentGPT.py}, and \texttt{agentClaude.py}.

\vs

On Windows, execute the following three commands one at a time for each API key. The example for \texttt{OPENAI\_API\_KEY} is as follows:

\begin{verbatim}
setx OPENAI_API_KEY "[sk-... API key]"
Exit
echo %OPENAI_API_KEY%
\end{verbatim}

In the code above, replace \texttt{[sk-... API key]} with your actual API key. The first command sets the environment variable permanently for your user account. The second command closes the Command Prompt to ensure the new environment variable is registered. The third command, run in a new Command Prompt window, prints the stored value to verify it was set correctly.

\vs

On Mac/Linux execute the following three commands one at a time for each API key. The example for \texttt{OPENAI\_API\_KEY} is as follows:

\begin{verbatim}
echo "export OPENAI_API_KEY='[sk-... API key]'" >> ~/.zshrc
source ~/.zshrc
echo $OPENAI_API_KEY
\end{verbatim}

In the code above, replace \texttt{'[sk-... API key]'} with your actual API key.

\vs

As reported by Bryan Fowler, if you experience errors trying to add your environmental variables in a Mac/Linux, the most likely problem is ownership of the \texttt{\~{}/.zshrc} file. In this case, execute

\begin{verbatim}
sudo chown $(whoami) ~/.zshrc
\end{verbatim}

And try to create the environmental variables again.

\vs

As reported by Drew Stephen regarding Mac, ``\textit{apparently it doesn't work on the default c shell but switching to the zsh lets it run}.'' You might need to close and reopen the terminal window from where you started for changes to become effective.

\section{The FOO Framework and Agent Feasibility}
\label{sec:foo}

\subsection{Minimum Agent Count}

The \textsc{Max} Multi-Agent Consensus System implements the FOO
algorithm for Byzantine-fault-tolerant consensus. Given that any
individual agent has an invalidation probability $p$, the minimum number
of agents required to achieve consensus confidence $C$ is:

\begin{equation}
N_{\text{foo}}(p, C)
  = \left\lceil \frac{\log(1 - C)}{\log(p)} \right\rceil
\label{eq:nfoo}
\end{equation}

Selected values are given in Table~\ref{tab:nfoo}.

\begin{table}[h]
\centering
\caption{Minimum agent count $N_{\text{foo}}$ by invalidation rate and
confidence target}
\label{tab:nfoo}
\begin{tabular}{ccc}
\toprule
Per-agent invalidation rate $p$ & Confidence $C$ & $N_{\text{foo}}$ \\
\midrule
0.30 & 0.95 & 4 \\
0.30 & 0.99 & 6 \\
0.20 & 0.95 & 3 \\
0.20 & 0.99 & 5 \\
0.10 & 0.99 & 3 \\
0.05 & 0.99 & 2 \\
\bottomrule
\end{tabular}
\end{table}

\subsection{Agent Feasibility Score}

The \emph{Agent Feasibility Score} (AFS) couples the hardware budget to
the FOO validity condition:

\begin{equation}
\text{AFS} = \frac{N_{\text{hw}}}{N_{\text{foo}}}
\label{eq:afs}
\end{equation}

A configuration is \emph{feasible} if and only if $\text{AFS} \geq 1.0$.
When $\text{AFS} < 1.0$, one or more of the following mitigations must
be applied: reduce the context window (increases $N_{\text{hw}}$), use a
smaller model (increases $N_{\text{hw}}$), or relax the confidence target
$C$ (reduces $N_{\text{foo}}$).

\subsection{Reference Configurations}

\paragraph{Dell VRAM-optimal (fast, no RAM spill).}
Five agents fitting entirely within 16\,GB VRAM, yielding
$\text{AFS} = 8.00$ at $p = 0.20$, $C = 0.99$:

\begin{center}
\begin{tabular}{lrr}
\toprule
Model & $M_{\text{agent}}$ (GB) & Location \\
\midrule
\texttt{gemma3:1b}    & 0.9 & VRAM \\
\texttt{phi4-mini:3.8b} & 2.8 & VRAM \\
\texttt{gemma3:4b}    & 2.8 & VRAM \\
\texttt{qwen3:4b}     & 2.8 & VRAM \\
\texttt{mistral:7b}   & 4.5 & VRAM \\
\midrule
Total & 13.8 & (of 16\,GB) \\
\bottomrule
\end{tabular}
\end{center}

\paragraph{Spark full-capability configuration.}
Mixed-tier configuration using 110\,GB unified memory:

\begin{center}
\begin{tabular}{llrr}
\toprule
Role & Model & $M_{\text{agent}}$ (GB) & AFS \\
\midrule
Orchestrator    & \texttt{qwen3:32b-ctx8k}            & 19.3 & 1.67 \\
Primary analyst & \texttt{llama3.3:70b-ctx8k}         & 42.5 & 0.67 \\
Critic A        & \texttt{mistral-small:22b}           & 14.0 & 2.33 \\
Critic B        & \texttt{phi4:14b}                   &  8.8 & 4.00 \\
Synthesiser     & \texttt{mistral-nemo:12b}            &  7.7 & 4.67 \\
\midrule
Total           &                                      & 92.3 & 0.67 \\
\bottomrule
\end{tabular}
\end{center}

\noindent The minimum AFS of 0.67 (from \texttt{llama3.3:70b} alone in
the pool) indicates that a full five-agent round cannot run simultaneously
on the Spark with this model selection. The operational solution is to
run the T1 models sequentially within a round and cache their outputs, or
to replace \texttt{llama3.3:70b-ctx8k} with \texttt{qwen3:32b-ctx8k} to
free memory for a fifth simultaneous agent.

\section{Single-Machine Configuration (Dell)}
\label{sec:single}

In single-machine mode the Dell runs both the \textsc{Alice} application
stack and the Ollama inference layer. This configuration is appropriate
for prompt development and unit-level agent testing where network
isolation is desirable.

\subsection{Environment Variable}

\begin{lstlisting}[style=bashstyle, caption={.env for single-machine mode}]
LLM_PROVIDER=ollama
LLM_BASE_URL=http://127.0.0.1:11434/v1
LLM_MODEL=mistral:7b
\end{lstlisting}

From WSL, the Windows Ollama instance is reachable via the WSL default
gateway rather than localhost:

\begin{lstlisting}[style=bashstyle, caption={Find the Windows gateway IP from WSL}]
ip route | grep default | awk '{print $3}'
# Typical result: 172.24.0.1 (varies by WSL version and config)
\end{lstlisting}

\begin{lstlisting}[style=bashstyle, caption={.env for WSL accessing Windows Ollama}]
LLM_PROVIDER=ollama
LLM_BASE_URL=http://172.24.0.1:11434/v1
LLM_MODEL=mistral:7b
\end{lstlisting}

\subsection{Model Selection for Single-Machine Development}

The Dell's 16\,GB VRAM constrains simultaneous model capacity. For
development work targeting the FOO framework, the following tier
assignment balances capability with VRAM residency:

\begin{itemize}
  \item \textbf{Primary agent:} \texttt{deepseek-r1:8b} (5.5\,GB),
        reasoning model with thinking mode; good proxy for larger
        consensus agents.
  \item \textbf{Fast critic:} \texttt{mistral-nemo:12b} (7.7\,GB),
        tools-capable, fast, NVIDIA-trained.
  \item \textbf{Utility agent:} \texttt{gemma3:1b} (0.9\,GB),
        sub-second responses for structured output tasks.
\end{itemize}

Total VRAM: 14.1\,GB out of 16\,GB. This leaves 1.9\,GB headroom for
the display server and background processes.

\section{Two-Machine Configuration}
\label{sec:two-machine}

In the two-machine setup the Dell hosts the application stack and the
Spark serves all LLM inference. The Spark's 110\,GB unified memory pool
allows multiple large models to remain resident simultaneously, which is
the hardware prerequisite for running a full FOO consensus round without
sequential eviction.

\subsection{Architecture}

\begin{verbatim}
+----------------------------+   HTTP/REST    +---------------------------+
|        Dell 7770           | <------------- |       NVIDIA Spark        |
|                            |   port 11434   |                           |
|  Alice FastAPI             |                |  Ollama (port 11434)      |
|  PostgreSQL 16             |                |  systemd: ollama.service  |
|  Agents                    |                |  OLLAMA_HOST=0.0.0.0      |
|  Python venv: ALICE        |                |  OLLAMA_NUM_CTX=8192      |
|                            |                |                           |
|  .env:                     |                |  Models resident:         |
|    LLM_BASE_URL=           |                |    qwen3:32b-ctx8k        |
|    http://10.2.14.24:      |                |    llama3.3:70b-ctx8k     |
|    11434/v1                |                |    mistral-nemo:12b       |
+----------------------------+                +---------------------------+
\end{verbatim}

\subsection{Environment Variables}

\begin{lstlisting}[style=bashstyle, caption={.env for two-machine mode (Dell application host)}]
# Two-machine mode: Spark serves all inference
LLM_PROVIDER=ollama
LLM_BASE_URL=http://10.2.14.24:11434/v1
LLM_API_KEY=not-required

# Per-agent model assignment (Max routing policy)
AGENT_MODEL_ORCHESTRATOR=qwen3:32b-ctx8k
AGENT_MODEL_ANALYST=llama3.3:70b-ctx8k
AGENT_MODEL_CRITIC=mistral-nemo:12b
AGENT_MODEL_FAST=mistral:7b
\end{lstlisting}

The \texttt{ProviderAdapter} protocol in \textsc{Max} makes switching
between single-machine and two-machine modes a one-line change to
\texttt{LLM\_BASE\_URL}. No application code changes are required.

\subsection{Connectivity Verification}

A connectivity test script (\texttt{connectivity\_test.py}) validates the
end-to-end path from the application host to the Ollama server. It tests
four conditions in sequence, stopping at the first failure:

\begin{enumerate}
  \item \textbf{TCP handshake:} confirms the host is reachable and the
        port is open.
  \item \textbf{MAC ARP check:} informational only; confirms whether the
        machines are on the same Layer~2 segment. Failure here is expected
        on a routed university network and does not indicate a problem.
  \item \textbf{Model list:} confirms Ollama is serving and returns the
        installed model list.
  \item \textbf{Inference:} sends a minimal chat completion request and
        validates the response.
\end{enumerate}

Machine profiles, ports, and test parameters are stored in
\texttt{connectivity\_config.json}, which is excluded from version control
(\texttt{.gitignore}) because it contains environment-specific IP and MAC
addresses. A \texttt{connectivity\_config.json.template} with placeholder
values is committed alongside the script.

\section{Capacity Testing}
\label{sec:capacity}

A capacity test script (\texttt{capacity\_test.py}) reads the model
registry from \\
\texttt{connectivity\_config.json}, computes the
theoretical $N_{\text{hw}}$ and AFS for each model, and empirically
validates each installed model by running a single inference call.

\subsection{Platform Considerations}

\begin{itemize}
  \item \textbf{ARM64 (Spark):} \texttt{qwen3:4b} fails with an
        \texttt{exit~status~2} crash in the Ollama ARM64 binary. This is
        an Ollama bug affecting this specific model and quantisation on
        ARM64. The same model works correctly on the x86\_64 Dell. The
        script detects architecture at runtime and skips ARM64-specific
        failures appropriately.

  \item \textbf{Modelfile variants:} Context-bounded variants (e.g.\
        \texttt{llama3.3:70b-ctx8k}) must be created separately on each
        machine. A model created via Modelfile on the Spark is not
        automatically available on the Dell. The script detects
        \emph{base installed but variant missing} and emits the exact
        shell commands to create the variant, using PowerShell syntax on
        Windows and bash heredoc syntax on Linux.

  \item \textbf{VRAM vs.\ RAM:} On the Dell, the script distinguishes
        models that fit entirely in VRAM (fast inference) from those that
        spill to system RAM (significantly slower). Two configuration
        recommendations are produced: a VRAM-optimal configuration
        (fastest, may not reach $N_{\text{foo}}$) and a RAM-extended
        configuration (reaches $N_{\text{foo}}$ if RAM permits, but with
        a performance warning for spilled models).
\end{itemize}

\subsection{Ollama CLI Reference}

\begin{lstlisting}[style=bashstyle, caption={Essential Ollama CLI commands}]
# List installed models
ollama list

# Show currently loaded models and memory usage
ollama ps

# Pull a model from the registry
ollama pull qwen3:32b

# Interactive chat session
ollama run mistral:7b

# One-shot prompt (no interactive session)
ollama run mistral:7b "Explain the rhizome model in one paragraph"

# Remove a model
ollama rm mistral:7b

# Test the OpenAI-compatible API directly
curl http://localhost:11434/v1/models
curl http://localhost:11434/v1/chat/completions \
  -H "Content-Type: application/json" \
  -d '{"model":"mistral:7b",
       "messages":[{"role":"user","content":"Hello"}]}'
\end{lstlisting}

\section{Deployment Parity}
\label{sec:parity}

A core principle of this infrastructure is that the development
environment mirrors production without containerisation. Both machines
run native Ubuntu (or WSL2 for the Windows host), Ollama is managed by
\texttt{systemd}, and configuration is injected via environment files
rather than code changes.

The two-tier configuration model is:

\begin{itemize}
  \item \textbf{Development:} \texttt{.env} file in the project root,
        loaded by \texttt{python-dotenv} with \texttt{override=False}.
  \item \textbf{Production:} \texttt{systemd} \texttt{EnvironmentFile=}
        directive pointing to a file outside the repository. No code
        changes required at deployment.
\end{itemize}

Switching from single-machine to two-machine mode, or from the Dell's
local Ollama to the Spark, requires only updating \texttt{LLM\_BASE\_URL}
in the environment file. The \texttt{ProviderAdapter} protocol
transparently routes requests to whichever endpoint is configured.

\section{Setting Up Ollama with a GUI on Ubuntu}

To transition from a command-line interface to a graphical experience, follow these steps to install the Ollama backend and the Open WebUI frontend.

\subsection{Step 1: Install Ollama Backend}
First, install the core Ollama service using the official installation script:
\begin{verbatim}
curl -fsSL https://ollama.com | sh
\end{verbatim}

\subsection{Step 2: Deploy Open WebUI via Docker}
Open WebUI is the most popular ChatGPT-like interface for Ollama. Run the following Docker command to pull the image and start the container:
\begin{verbatim}
docker run -d -p 3000:8080 \
  --add-host=host.docker.internal:host-gateway \
  -v open-webui:/app/backend/data \
  --name open-webui \
  ghcr.io/open-webui/open-webui:main
\end{verbatim}

\subsection{Step 3: Access the Interface}
Once the container is running, open your preferred web browser and navigate to:
\begin{quote}
    \texttt{http://localhost:3000}
\end{quote}
You will be prompted to create an admin account (stored locally) to begin chatting with your models.

\subsection{Technical Considerations}
\begin{itemize}
    \item \textbf{GPU Support:} For NVIDIA users, ensure the \texttt{nvidia-container-toolkit} is installed to enable hardware acceleration within Docker.
    \item \textbf{External Access:} To access this GUI from another machine, set the environment variable \texttt{OLLAMA\_HOST=0.0.0.0} in your systemd service configuration.
\end{itemize}

